% !TEX program = xelatex
\documentclass[a4paper,11pt]{article}

% ============================================================
%  Packages
% ============================================================
\usepackage[UTF8]{ctex}
\usepackage[margin=2.5cm, top=3cm, bottom=3cm]{geometry}
\usepackage{graphicx}
\usepackage{xcolor}
\usepackage{hyperref}
\usepackage{booktabs}
\usepackage{longtable}
\usepackage{array}
\usepackage{tabularx}
\usepackage{fancyhdr}
\usepackage{titlesec}
\usepackage{enumitem}
\usepackage{float}
\usepackage{multicol}
\usepackage{tikz}

% ============================================================
%  Colors
% ============================================================
\definecolor{brandblue}{RGB}{0, 51, 102}
\definecolor{accentblue}{RGB}{0, 102, 179}
\definecolor{lightgray}{RGB}{240, 240, 240}
\definecolor{darkgray}{RGB}{80, 80, 80}
\definecolor{rulergray}{RGB}{180, 180, 180}

% ============================================================
%  Hyperlinks
% ============================================================
\hypersetup{
  colorlinks=true,
  linkcolor=brandblue,
  urlcolor=accentblue,
  pdfauthor={Shanghai Qiongpei Technology Co., Ltd.},
  pdftitle={Statement of Technical Capabilities},
}

% ============================================================
%  Header / Footer
% ============================================================
\pagestyle{fancy}
\fancyhf{}
\fancyhead[L]{\small\color{darkgray} Shanghai Qiongpei Technology Co., Ltd.}
\fancyhead[R]{\small\color{darkgray} Confidential}
\fancyfoot[C]{\small\thepage}
\fancyfoot[R]{\small\color{darkgray} April 2026}
\renewcommand{\headrulewidth}{0.5pt}
\renewcommand{\footrulewidth}{0.3pt}
\renewcommand{\headrule}{\hbox to\headwidth{\color{brandblue}\leaders\hrule height \headrulewidth\hfill}}
\renewcommand{\footrule}{\hbox to\headwidth{\color{rulergray}\leaders\hrule height \footrulewidth\hfill}}

% ============================================================
%  Section Styling
% ============================================================
\titleformat{\section}
  {\normalfont\Large\bfseries\color{brandblue}}
  {\thesection.}{0.5em}{}
  [\vspace{-0.5em}{\color{brandblue}\rule{\textwidth}{0.5pt}}]

\titleformat{\subsection}
  {\normalfont\large\bfseries\color{accentblue}}
  {\thesubsection}{0.5em}{}

% ============================================================
%  Custom environments
% ============================================================
\newcommand{\capbox}[2]{%
  \begin{center}
  \fcolorbox{accentblue}{accentblue!5}{%
    \parbox{0.92\textwidth}{\textbf{#1}\vspace{0.3em}\newline #2}%
  }
  \end{center}
}

\setlist[itemize]{leftmargin=1.5em, itemsep=0.3em}
\setlist[enumerate]{leftmargin=1.5em, itemsep=0.3em}

% ============================================================
%  Document
% ============================================================
\begin{document}

% ────────────────────────────────────────────────────────────
%  COVER PAGE
% ────────────────────────────────────────────────────────────
\thispagestyle{empty}
\begin{center}

\vspace*{4cm}

{\color{brandblue}\rule{\textwidth}{1.5pt}}

\vspace{1.5cm}

{\Huge\bfseries\color{brandblue} Statement of\\[0.4em] Technical Capabilities}

\vspace{0.8cm}

{\Large\color{darkgray} in Electromagnetics and\\[0.3em] Electromechanical Control}

\vspace{1.5cm}

{\color{brandblue}\rule{\textwidth}{1.5pt}}

\vspace{2cm}

{\large Prepared for the\\[0.3em]
\textbf{Immersive Hiking Simulator Treadmill Project}}

\vspace{3cm}

{\Large\bfseries Shanghai Qiongpei Technology Co., Ltd.}\\[0.5em]
{\large 上海穹沛科技有限公司}

\vspace{1cm}

{\large April 2026}\\[0.5em]
{\color{darkgray} Document Classification: Confidential}

\end{center}

\newpage

% ────────────────────────────────────────────────────────────
%  TABLE OF CONTENTS
% ────────────────────────────────────────────────────────────
\tableofcontents
\newpage

% ────────────────────────────────────────────────────────────
%  I. EXECUTIVE SUMMARY
% ────────────────────────────────────────────────────────────
\section{Executive Summary}

This document presents the technical capabilities of Shanghai Qiongpei Technology Co., Ltd. as they relate to the Immersive Hiking Simulator Treadmill Project. Our assessment is grounded in five production-validated technology platforms that directly address the core engineering challenges of this project.

\capbox{Core Argument}{%
The central challenge of the Immersive Hiking Simulator Treadmill is not the selection of any single component. It is the \textbf{system-level integration} of electromagnetic arrays, drive control, multi-sensor feedback, haptic response, audio-visual synchronization, and safety mechanisms into a reliable, maintainable engineering system.\\[0.5em]
This class of problem---multi-actuator, multi-sensor, real-time closed-loop control with stringent safety requirements---is precisely what our team has been building and shipping for the past three years.}

\vspace{0.5em}

\noindent Our existing technology capabilities map directly to the project requirements:

\begin{center}
\begin{tabularx}{\textwidth}{>{\bfseries}l X >{\itshape}X}
\toprule
\textbf{Capability Area} & \textbf{What We Have Built} & \textbf{Project Relevance} \\
\midrule
Motor Drive \& Control & Self-developed FOC drive boards, current regulation, thermal protection & Electromagnetic array drive and protection \\
Multi-Actuator Systems & 16-DOF robotic platform with multi-sensor fusion & Multi-channel array coordination \\
Real-Time Control & FSM architecture, gait control, dynamic state switching & Terrain-adaptive feedback control \\
Safety \& Fault Handling & E-Stop, safety gates, fault recovery, real-time telemetry & System safety and fault tolerance \\
\bottomrule
\end{tabularx}
\end{center}

% ────────────────────────────────────────────────────────────
%  II. TECHNICAL CAPABILITIES
% ────────────────────────────────────────────────────────────
\section{Existing Technical Capabilities}

\subsection{Electromagnetic Drive and Low-Level Control}

Our self-developed motor drive platform provides production-validated capabilities in electromagnetic actuation and precision current control.

\vspace{0.5em}
\noindent\textbf{Hardware Platform:}
\begin{itemize}
  \item High-performance Field-Oriented Control (FOC) on STM32G4 microcontroller
  \item CAN bus communication interface for multi-device coordination
  \item Triple closed-loop control: position, velocity, and torque
  \item High-precision current sampling with offset calibration
  \item Hardware-level overtemperature protection and fault shutdown logic
\end{itemize}

\vspace{0.5em}
\noindent\textbf{Demonstrated Competencies:}
\begin{itemize}
  \item Electromagnetic actuator drive circuit design and validation
  \item Precision current regulation and power envelope management
  \item High-frequency PWM control with sub-millisecond response latency
  \item Thermal modeling, protection threshold definition, and derating strategy
  \item Coordinated design across drive boards, control firmware, and host systems
\end{itemize}

\capbox{Relevance to Electromagnetic Array Control}{%
Although the project employs electromagnetic arrays rather than rotary motors, the two are closely analogous at the drive-and-control level. Both involve coil excitation waveform design, precision current regulation, high-frequency switching response, thermal boundary management, and fault protection. Our drive platform addresses all of these requirements with proven, production-grade firmware.}

\subsection{Multi-Actuator, Multi-Sensor Integrated Systems}

Our industrial wheeled-legged quadruped robot platform is a multi-actuator, multi-sensor, real-time closed-loop control system that has achieved stable locomotion across six terrain types.

\vspace{0.5em}
\noindent\textbf{System Specifications:}
\begin{itemize}
  \item Wheeled-legged quadruped configuration with 16 degrees of freedom
  \item Self-developed motor drive and actuation system
  \item Multi-modal sensor suite: LiDAR, depth cameras, IMU, joint encoders
  \item Unified control architecture spanning motion control, state machines, safety logic, and supervisory interfaces
  \item Demonstrated stable traversal across six terrain types (flat, slope, stairs, gravel, grass, uneven ground)
\end{itemize}

\vspace{0.5em}
\noindent\textbf{Demonstrated Competencies:}
\begin{itemize}
  \item Parallel coordination of 16+ actuators with synchronized timing
  \item Multi-sensor data fusion with deterministic timing guarantees
  \item Real-time closed-loop control at 50\,Hz with sub-2\,ms inference latency
  \item System-level integration, commissioning, and stability validation
  \item Terrain-differentiated control strategy design and switching
\end{itemize}

\capbox{Relevance to Multi-Channel Array Control}{%
The electromagnetic treadmill array is fundamentally a multi-actuator system requiring spatially distributed, individually addressable control---the same class of problem as coordinating 16 independent motor drives on a quadruped platform. Our demonstrated capability to maintain stable, synchronized control across multiple actuators under dynamic conditions transfers directly to array-level electromagnetic control.}

\subsection{Real-Time Motion Control and State Machine Architecture}

Our motion control stack provides a production-grade framework for real-time state management and dynamic behavior switching.

\vspace{0.5em}
\noindent\textbf{Architecture:}
\begin{itemize}
  \item Finite State Machine (FSM) with formally defined states and safe transitions
  \item Remote command interface with priority-based control arbitration
  \item 50\,Hz control loop with real-time inference and PD gain scheduling
  \item Modular architecture supporting third-party integration via standard APIs
\end{itemize}

\vspace{0.5em}
\noindent\textbf{Demonstrated Competencies:}
\begin{itemize}
  \item Design and implementation of multi-state control logic with safe transitions
  \item Real-time dynamic parameter switching (PD gains, motion profiles, terrain modes)
  \item Coordination between control dispatch, command arbitration, and safety interlocks
  \item SDK-level API design for system integration by third-party developers
\end{itemize}

\capbox{Relevance to Dynamic Terrain Feedback}{%
Electromagnetic array control for a hiking simulator is not a static output problem. It is a dynamic system that must continuously adapt to footstep events, walking speed, terrain mode selection, and safety status. This is precisely the class of problem that state machines and real-time control frameworks are designed to solve. Our control architecture can be directly adapted to manage terrain-mode switching, cadence-responsive feedback, and multi-modal synchronization.}

\subsection{Safety Control, Fault Handling, and Remote Maintenance}

Our runtime safety framework provides comprehensive protection and fault management for electromechanical systems.

\vspace{0.5em}
\noindent\textbf{Safety Architecture:}
\begin{itemize}
  \item Hardware E-Stop safety layer integrated into the system architecture
  \item Safety gates: every actuation command passes through capability checks before execution
  \item Structured fault detection, classification, and recovery logic
  \item Real-time telemetry for all critical system parameters
  \item End-to-end auditability with timestamped event logs
\end{itemize}

\vspace{0.5em}
\noindent\textbf{Direct Correspondence to Project Requirements:}

\begin{center}
\begin{tabularx}{\textwidth}{X X}
\toprule
\textbf{Project Safety Requirement} & \textbf{Existing Capability} \\
\midrule
Electromagnetic array anomaly protection & Hardware fault detection + automatic shutdown \\
Overtemperature protection & Thermal monitoring + derating + E-Stop \\
Emergency stop logic & Hardware E-Stop + software safety gates \\
Sensor fault handling & Structured fault classification + recovery \\
Long-duration operational stability & Telemetry + thermal management + watchdogs \\
Ongoing firmware maintenance & Remote update capability with rollback support \\
\bottomrule
\end{tabularx}
\end{center}

% ────────────────────────────────────────────────────────────
%  III. PROJECT REQUIREMENT MAPPING
% ────────────────────────────────────────────────────────────
\section{Direct Mapping to Project Requirements}

The Immersive Hiking Simulator Treadmill Project presents six core technical challenges. Each maps to specific, validated capabilities within our technology stack.

\begin{center}
\begin{tabularx}{\textwidth}{>{\raggedright\arraybackslash}p{0.55\textwidth} >{\raggedright\arraybackslash}p{0.38\textwidth}}
\toprule
\textbf{Project Challenge} & \textbf{Mapped Capability} \\
\midrule
\textbf{1. Stable, adjustable resistance feedback} from electromagnetic array and copper woven conductive layer
& Precision current regulation, closed-loop torque control, thermal protection \\
\addlinespace
\textbf{2. Spatially distributed multi-channel control} rather than single constant output
& Parallel multi-actuator coordination with independent per-channel control \\
\addlinespace
\textbf{3. Dynamic feedback triggered} by footsteps, speed variation, and terrain mode
& FSM-based state switching, real-time event-driven control, cadence detection \\
\addlinespace
\textbf{4. Synchronization} of electromagnetic feedback with haptics, audio, video, and safety
& Deterministic timing architecture, multi-subsystem coordination, priority arbitration \\
\addlinespace
\textbf{5. Long-duration stability} with controlled temperature rise and abnormal-condition handling
& Thermal management framework + safety layer + real-time telemetry monitoring \\
\addlinespace
\textbf{6. Simulation-driven design convergence} for magnetic field distribution, coupling, and thermal boundaries
& Electromagnetic simulation methodology + parametric validation workflow \\
\bottomrule
\end{tabularx}
\end{center}

\vspace{1em}

\capbox{Engineering Match, Not Conceptual Match}{%
The correspondence between our capabilities and this project's requirements is not theoretical. It is grounded in \textbf{production-validated} systems: drives that have powered real actuators, control stacks that have maintained stability across real terrain, safety systems that have protected real hardware, and maintenance systems that have updated real devices in the field.}

% ────────────────────────────────────────────────────────────
%  IV. ASSESSMENT OF PROJECT CHALLENGES
% ────────────────────────────────────────────────────────────
\section{Professional Assessment of Project Challenges}

From an engineering-implementation standpoint, the most difficult aspects of this project are not in the electromagnetic array itself, but in the following system-level integration challenges:

\begin{enumerate}
  \item \textbf{Subjective feel calibration.} Tuning the damping and resistance feedback to feel natural and immersive requires iterative testing with human subjects. This cannot be solved through simulation alone---it requires rapid prototyping capability and a control system that supports real-time parameter adjustment.

  \item \textbf{Terrain differentiation fidelity.} Different terrain types (rock trail, forest path, sand, snow) must produce perceptibly distinct feedback profiles. This requires a control architecture that supports multiple terrain modes with smooth transitions between them.

  \item \textbf{Cadence consistency.} The feedback must remain consistent under varying left-right foot triggering patterns and walking speeds. This is a multi-sensor fusion and timing-coordination problem.

  \item \textbf{Multi-modal synchronization.} Electromagnetic resistance feedback must be precisely synchronized with haptic actuators, spatial audio, and visual projection. Latency mismatches as small as 50\,ms degrade the immersive experience.

  \item \textbf{Thermal management under sustained operation.} Commercial installations may operate 8--12 hours per day. Thermal management and stability must be engineered from the ground up, not added as an afterthought.

  \item \textbf{Retrofit-space constraints.} Integrating electromagnetic, mechanical, electrical, and control systems within the physical constraints of a treadmill form factor requires coordinated multi-discipline design.
\end{enumerate}

\vspace{0.5em}

The common characteristic of these challenges is that they must be resolved through \textbf{progressive validation and de-risking}---combining electromechanical control experience, real-time systems experience, and system-level integration capability. Theoretical analysis alone is insufficient. Our team's capabilities have been built around precisely this class of system-level engineering problem.

% ────────────────────────────────────────────────────────────
%  VI. CONCLUSION
% ────────────────────────────────────────────────────────────
\section{Conclusion}

Shanghai Qiongpei Technology Co., Ltd. possesses the following tangible, production-validated technical capabilities that are directly relevant to the Immersive Hiking Simulator Treadmill Project:

\begin{enumerate}
  \item \textbf{Self-developed FOC drive and low-level electromagnetic control} --- precision current regulation, triple closed-loop control, and thermal protection on custom hardware
  \item \textbf{Multi-actuator, multi-sensor robotic system development} --- 16-DOF parallel actuation, multi-modal sensing, and terrain-adaptive control
  \item \textbf{Real-time motion control and state machine architecture} --- FSM-based control, dynamic parameter switching, and SDK-level integration interface
  \item \textbf{Safety control and fault handling} --- E-Stop safety layer, safety gates, structured fault recovery, and real-time telemetry
  \item \textbf{Electromagnetic simulation and parametric design convergence} --- Methodology for magnetic-field distribution analysis, thermal-boundary verification, and prototype parameter definition
  \item \textbf{System-level integration and commissioning} --- Demonstrated ability to bring multi-discipline systems from prototype to stable operation
\end{enumerate}

\vspace{1em}

\capbox{Summary}{%
On the basis of these capabilities, our team is well positioned to undertake the electromagnetic system design, drive control development, sensing closed-loop implementation, system integration, and commissioning and validation work required for the prototype phase of this project.\\[0.5em]
We welcome the opportunity to discuss the project in further detail and to define a concrete engagement plan.}

\vspace{2cm}

\begin{flushright}
\textbf{Shanghai Qiongpei Technology Co., Ltd.}\\
上海穹沛科技有限公司\\[1em]
April 2026
\end{flushright}

\end{document}
